% =============================================================
% THE BALTIC DATA GAP AND PREVIOUS ESTIMATES
% =============================================================
\section{The Baltic data gap and previous estimates}
\label{sec:gap}

The absence of the Baltic states from the quantitative economic history of Europe is a consequence of their history rather than of neglect. Latvia and Estonia have a remarkably brief record of state continuity, so that little impetus existed before 1918 to maintain the consistent statistical records that state-building elsewhere produced; the Russian Empire into which all three future republics had been absorbed in the eighteenth century was itself a country of weak statistical coverage \parencite[219--20]{gregrussianempire}. Regular statistical publication began with independence, and Latvia was among the first ten states in the world whose statistical office published estimates of national income, for the mid-1930s \parencite{norkuslatvia1935}. But the period from 1940 to 1945 saw the country occupied three times, with changes of borders, economic regime and statistical practice and the physical destruction of part of the archival record \parencite[140]{dagys1959}, and after 1945 the Soviet authorities discouraged study of the capitalist decades \parencite[244]{morysivanovinterwarinmorys2021}. For the Soviet period the difficulty is of a different kind. The republics' `national income' statistics, compiled from the 1950s, measured net material product, which excludes most services and is net of depreciation; the ratio of net material product to GDP varies across countries and over time, so that conversion is unsafe \parencites[581--93]{ceska1989}{Campbell1985}; and the series were compiled at `comparable prices' that understated inflation and overstated growth, a practice that Adlys, deputy head of the Lithuanian statistical committee, acknowledged in a published interview of 1988: five-year plans `were not fulfilled but we showed that they were' \parencite{adlys1988}. The official index of Latvian industrial production at constant prices rises between 1940 and 1990 by several times the increase recorded by the physical-product index built in the next section.

Neither the Baltic statistical offices nor Eurostat, the IMF or the World Bank publish GDP for the Baltic states before 1990, and the World Bank's own reconstructions of 1980--91, made by its missions of 1992--3, were discontinued \parencites{wbtransitionlatv}[32--3]{worldtables1995}. The Maddison Project Database reports Latvian GDP per capita for 1973 and annually from 1980. Its 1990 level descends from Maddison's application of Bolotin's republic-to-Union ratios, taken for 1991 rather than 1990, to a Soviet total derived from a bilateral Austria--USSR comparison; its earlier years are extrapolated backwards with net material product growth adjusted by the Union-wide ratio of GDP to net material product, a short cut that Maddison himself regarded as inadmissible for the Soviet Union as a whole \parencites[307--8]{Maddison1998b}[181--2]{maddisonworldeconomy}. The resulting series shows Latvian income per head rising by 5 per cent in 1989, when every physical indicator was already falling, and stands \mpdgap{1990} per cent above Bolotin's independent benchmark for 1990 (Table~\ref{tab:previous}). The only republic-level estimates in the Maddison tradition are the indirect approximations of \textcite{norkusmark2021} for 1913, 1922, 1929 and 1938, obtained from real wages and the structure of employment, and the figures of \textcite[187--90]{roseswolf2010} for the same interwar years, which rest on an undocumented procedure and make interwar Latvia Europe's fastest-growing economy and, by 1938, richer than Finland.

For Latvia the direct estimates are recent and confined to the interwar years. \textcite{norkuslatvia1935} produced an output-side benchmark for 1935 within the 2008 System of National Accounts, with gross value-added for twenty industries at basic and market prices and a conversion into international dollars by a bilateral price comparison with Sweden; \textcite{Klimantas2024} extrapolated it to annual series for 1920--39 from sectoral volume indices, resolving the long dispute between Aizsilnieks, for whom national income per head did not grow under the Ulmanis dictatorship of 1934--40, and the Latvian historians who regarded those years as prosperous: growth continued after 1934, but Latvia was not Europe's fastest-growing economy in 1929--38, and its re-industrialisation was mainly import-substituting \parencite{latvija19141945vesture}. The series is the anchor of the present paper. For the Soviet period there is nothing comparable. The Western Sovietologists who studied the Latvian economy worked from official Soviet aggregates and could at best qualify them: King's study of economic policy in occupied Latvia and Viksnins's evaluations of Latvian growth under Brezhnev and Gorbachev document the structure of investment, labour and consumption but offer no measure of real output comparable across time or countries \parencite{king_1965_economic,viksnins_1981_evaluating,viksnins_1986_latvian}, and Koropeckyj's attempt to place the national income of the Baltic republics in 1970 at British and Australian levels illustrates the hazards of such comparisons \parencite[72]{bs_livingstandards1976}. Latvian historiography since 1991 has produced a rich institutional account of the Soviet economy, culminating in the collective economic history edited by one of the present authors and in the reassessment of Latvia's fiscal relations with Moscow, which overturned the Soviet claim that the republic was a net recipient of Union investment \parencite{gatis_lv_eh,gatis_compensations}; but it, too, has had to rely on Soviet output indices whose growth is several times the real one. The Lithuanian and Estonian series now published close two-thirds of the gap \parencite{klimantas2026soviet,klimantas2025estonia}. The Latvian series presented in the next section closes the rest, and allows the three Baltic economies, which are compared with each other more often than with anyone else, to be compared on a common footing for the first time.
